% T18 -- Circuits \'electroniques, composants (diode, transistor) -- banque d'exercices
% Sources (v\'erifi\'ees dans le texte extrait des PDF originaux) :
%   - Exercice 1 : E2_Disciplinaire_IE_2022, questions 26 ET 27 (module TRT200, surveillance des
%     accumulateurs). \'Epreuve intitul\'ee EDE ENE 1 "\'epreuve \'ecrite disciplinaire" (et NON
%     "appliqu\'ee" : l'\'epreuve appliqu\'ee commune 2022 porte sur le banc \`a rouleaux AutoScan Fi).
%     Donn\'ees du sujet : tension de seuil des diodes 0,6 V et V_ref du CAN = 3,5 V
%     (E2_Disciplinaire_IE_2022.md l.453 et l.455) ; \'enonc\'es Q26 et Q27 l.473 et l.475.
%     LOT 3 -- CET EXERCICE A \'ET\'E ENTI\`EREMENT R\'E\'ECRIT. L'ancienne version (protection contre
%     une SURtension, \'ecr\^etage \`a V_ref + 0,6 = 4,1 V) \'etait un doublon exact de la QT3 de la
%     fiche T18 ET reprenait son interpr\'etation physique, laquelle est CONTREDITE par le corrig\'e
%     officiel : Rapport_Jury_2022.md l.343 -- "Lorsque les accumulateurs sont en charge, le courant
%     de charge circule de la borne + vers la borne - des accumulateurs. La diode est ainsi bloqu\'ee.
%     Lorsque les accumulateurs d\'ebitent vers le syst\`eme le courant change de sens induisant une
%     tension n\'egative qui peut \^etre nuisible pour l'entr\'ee du CAN. La diode devient ainsi
%     passante (VAK > 0,6V). La tension ainsi appliqu\'ee sur l'entr\'ee du CAN [...] ne pourra
%     d\'epasser -0,6V." La protection est donc un \'ecr\^etage \`a -0,6 V (tension N\'EGATIVE), pas
%     \`a +4,1 V. R\^ole de R0 : Rapport_Jury_2022.md l.339 -- "R0 est une r\'esistance de shunt. On
%     r\'ecup\`ere une tension image du courant de charge."
%   - Exercice 2 : E1_Sujet_2013_commun_SII, question 48 ("\'epreuve de synth\`ese de sciences
%     industrielles", commune aux options) : r\^ole de l'onduleur photovolta\"ique, 24 onduleurs
%     EI3300. La r\'eponse attendue (commande MLI des transistors) est celle du corrig\'e officiel
%     E1_Corrige_2013_commun_SII question 48 ; elle a \'et\'e retir\'ee de l'\'enonc\'e, o\`u elle \'etait
%     donn\'ee, pour redevenir une v\'eritable question
%   - Exercice 3 (droite de charge diode) : exercice in\'edit, caract\'eristique statique
%     g\'en\'erique (non issue d'un composant commercial pr\'ecis)
%   - Exercice 4 (hacheur, viabilit\'e des \'etats $Q_1$/$D_1$, lois de Kirchhoff) : exercice IN\'EDIT
%     inspir\'e du convertisseur CV3-1 du sujet 2025 (E2_Disciplinaire_IE_2025, DT10, Q24 \`a Q32).
%     Dans ce sujet, CV3-1 est un convertisseur \'EL\'EVATEUR (point de repos
%     $E_{b0}/V_{DC0} = 1-\alpha_0 = 48/230$) ; l'exercice traite ici la structure duale, d\'evoltrice.
%     La r\'ef\'erence "questions 24 \`a 27" a \'et\'e retir\'ee du titre au profit d'une mention
%     "inspir\'e de", et une remarque du corrig\'e explicite le passage \`a la relation \'el\'evatrice.
%     LOT 3 -- la question 2 demandait pr\'ec\'edemment les \'equations diff\'erentielles de $i_{L1}$ sur
%     $]0;\alpha T]$ et $]\alpha T;T]$, c'est-\`a-dire mot pour mot le travail de la QT1 de la fiche
%     (2025 Q27-Q28). Elle a \'et\'e remplac\'ee par les lois g\'en\'erales de Kirchhoff (2025 Q26,
%     E2_Disciplinaire_IE_2025.md l.312), non trait\'ees par la fiche, et le sch\'ema a re\c cu le
%     fl\'echage complet de Q25 ($v_{Q1}$, $v_{D1}$, $i_{Q1}$, $i_{D1}$).
%   - Exercice 5 (transistor bipolaire en commutation, pilotage relais) : exercice IN\'EDIT.
%     LOT 3 -- converti en exercice de LECTURE DE DOCUMENT TECHNIQUE (seul docTech de la notion).
%     L'extrait de datasheet est celui du BC337-25 (Fairchild Semiconductor, "NPN General Purpose
%     Amplifier", r\'ef. document 33716-25 Rev. B, (c) 1997), r\'ecup\'er\'e et v\'erifi\'e par WebSearch
%     puis lecture du PDF constructeur (datasheet.octopart.com/BC337-Fairchild-Semiconductor-
%     datasheet-10379962.pdf), recoup\'e avec la version Fairchild BC337/338 Rev. A2 (ao\^ut 2002).
%     TOUTES les valeurs de l'encadr\'e docTech sont des valeurs constructeur relev\'ees, AUCUNE
%     n'est invent\'ee : V_CEO=45 V, V_CES=50 V, V_EBO=5,0 V, I_C continu=1,0 A, P_D=625 mW,
%     T_J/T_stg=-55..+150 degC ; h_FE 160..400 (V_CE=1,0 V, I_C=100 mA) et h_FE min 40
%     (V_CE=1,0 V, I_C=500 mA) ; V_CE(sat) max 0,7 V (I_C=500 mA, I_B=50 mA) ; V_BE(on) max 1,2 V
%     (V_CE=1,0 V, I_C=500 mA). Les valeurs d'application (V_cc=12 V, R_relais=100 ohm, V_cmd=5 V)
%     restent celles, in\'edites, de l'\'enonc\'e.
%   - Exercice 6 (thyristor, pont P1) : formule $\langle u_S\rangle = \frac{V_m}{2\pi}(1+\cos\alpha)$
%     red\'emontr\'ee ici -- corrige un facteur 2 manquant dans fiches/T18_contenu.md (la forme qui y
%     figure, $V_m(1+\cos\alpha)/\pi$, ne redonne pas la valeur moyenne usuelle $V_m/\pi$ du
%     redressement demi-onde non command\'e $\alpha=0$, ce qui la rend incoh\'erente)
% Ordre des exercices : gradation croissante 1-1-2-2-2-3. L'exercice photovolta\"ique (niveau 1) a
% \'et\'e remont\'e en 2e position ; il figurait en 4e, apr\`es deux exercices de niveau 2.
% Sch\'emas (contr\^ole visuel sur rendu PNG) : ex. 4 -- $D_1$ retourn\'ee (cathode au n\oe ud de
%   commutation, anode au rail 0 V) : dessin\'ee en sens inverse, elle court-circuitait $V_{DC}$ \`a
%   l'amor\c cage de $Q_1$ et contredisait le corrig\'e ; fl\'echage des tensions $V_{DC}$, $u_{L1}$,
%   $E_b$ ajout\'e (questions 2 et 3 exigent une convention explicite). Ex. 3 -- caract\'eristique de
%   la diode recal\'ee (coude 0,54 V au lieu de 0,62 V) pour que l'intersection avec la droite de
%   charge tombe r\'eellement sur le point $Q$ annonc\'e (0,66 V ; 20 mA) : le point trac\'e \'etait
%   auparavant hors de la courbe.
% Sens des diodes (contr\^ole LOT 3, les 4 diodes du fichier) : ex.1 $D_1$ anode au 0 V / cathode
%   \`a l'entr\'ee ADC1 (\'ecr\^etage des tensions N\'EGATIVES, conforme au corrig\'e jury 2022) ;
%   ex.3 $D$ anode en haut / cathode vers la masse (sens direct du courant de la maille) ;
%   ex.4 $D_1$ cathode au n\oe ud de commutation (roue libre d'un Buck, INCHANG\'E depuis le lot 2) ;
%   ex.5 $D$ anode au collecteur / cathode \`a $V_{cc}$ (roue libre, INCHANG\'E depuis le lot 2).
\expandafter\ifx\csname CAPETINCLUDEMODE\endcsname\relax
\documentclass[11pt,a4paper]{article}
% preamble.tex — packages et commandes communes a toutes les fiches CAPET SII IE
\input{glyphtounicode}
\pdfgentounicode=1
\usepackage{cmap}
\usepackage[T1]{fontenc}
\usepackage[utf8]{inputenc}
\usepackage[french]{babel}
% babel[french] active a \begin{document} des raccourcis de ponctuation (:;!?)
% qui cassent la syntaxe to[...] de circuitikz (erreur "+ or - expected") ;
% on les desactive, ce qui ne change que l'espacement fin avant ces signes.
\AtBeginDocument{\shorthandoff{:;!?}}
\usepackage[a4paper, top=1.6cm, bottom=1.6cm, left=1.4cm, right=1.4cm]{geometry}
\setlength{\headheight}{14.5pt}
\usepackage{amsmath, amssymb, amsfonts}
\usepackage{siunitx}
\sisetup{locale=FR, output-decimal-marker={,}, inter-unit-product=\ensuremath{{}\cdot{}},
         group-minimum-digits=4}
% Unites non fournies par siunitx, declarees ici une fois pour toutes (fiches + exercices).
% Les gardes \ifdefined evitent tout conflit avec une redeclaration locale restee dans un
% fichier : ne pas retirer les gardes, ne pas redeclarer ces unites dans les fichiers.
\DeclareSIUnit{\jour}{j}
\ifdefined\var\else\DeclareSIUnit{\var}{var}\fi   % puissance reactive
\ifdefined\vad\else\DeclareSIUnit{\vad}{VAD}\fi   % puissance deformante
\ifdefined\tr\else\DeclareSIUnit{\tr}{tr}\fi      % tour (tr/min)
\usepackage{array, booktabs, multirow, tabularx}
\usepackage{xcolor}
\usepackage[most, breakable]{tcolorbox}
\usepackage{multicol}
\usepackage{enumitem}
\usepackage{mathtools}
\usepackage{lmodern}
% Refonte visuelle 2026 (direction "moderne coloré") : Lato pour le texte,
% maths sans-serif assorties via newtxsf (charge apres amsmath/amssymb).
\usepackage[default]{lato}
\usepackage{newtxsf}
\usepackage{microtype}
\usepackage{fancyhdr}
\usepackage{tikz}
\usetikzlibrary{positioning,arrows.meta}
\usepackage{circuitikz}
\usepackage{lastpage}

% ── Couleurs par domaine (palette modernisée 2026) ────────────────────
\definecolor{colElec}{RGB}{29, 78, 216}        % Électrotechnique — bleu vif
\definecolor{colElecLight}{RGB}{219, 234, 254} % fond bleu clair
\definecolor{colAuto}{RGB}{21, 128, 61}        % Automatique — vert
\definecolor{colMeca}{RGB}{194, 65, 12}        % Mécanique — orange
\definecolor{colTherm}{RGB}{185, 28, 28}       % Thermique — rouge
\definecolor{colQT}{RGB}{232, 244, 233}        % (hérité, plus utilisé par questiontype)
\definecolor{colJury}{RGB}{255, 248, 225}      % fond jury (crème)
\definecolor{colJuryFrame}{RGB}{202, 138, 4}   % bordure/bandeau jury (ambre)
\definecolor{colGrayLight}{RGB}{245, 245, 247} % fond gris neutre
\definecolor{colSeparator}{RGB}{180, 180, 180} % gris séparateur
\definecolor{colExo}{RGB}{8, 145, 178}          % Banque d'exercices — cyan (distinct des fiches)
% ── Encre et gris de la refonte ───────────────────────────────────────
\definecolor{colInk}{RGB}{23, 32, 42}          % texte courant (quasi-noir)
\definecolor{colSlate}{RGB}{83, 98, 118}       % gris ardoise (méta, en-têtes)
\definecolor{colRule}{RGB}{203, 213, 225}      % filets clairs
\definecolor{colJuryBg}{RGB}{255, 248, 225}    % fond crème jury
\definecolor{colPresence}{RGB}{83, 98, 118}    % bloc 5 (présence dans les sujets)
\AtBeginDocument{\color{colInk}}

% ── Style encadré exercice + corrigé (banque d'exercices) ─────────────
\tcbset{
  exoBox/.style={
    breakable,
    enhanced,
    colback=colExo!8!white,
    colframe=colExo,
    colbacktitle=colExo,
    coltitle=white,
    fonttitle=\bfseries\small,
    left=4pt, right=4pt, top=3pt, bottom=3pt,
    before skip=5pt, after skip=5pt,
  }
}

% ── Bascule enonce / corrige (banque d'exercices) ─────────────────────
% \CAPETCORRIGE est defini en ligne de commande (pdflatex "\def\CAPETCORRIGE{1}\input{...}")
% pour compiler la variante corriges d'un meme fichier source ; par defaut (non defini),
% la compilation produit la variante enonces. \CAPETSHOW{enonce}{corrige} choisit la branche.
%
% TROIS MODES (refonte 2026-09-20, manuel de revision) :
%   - aucun drapeau   -> imprime l'enonce            (variante enonces, INCHANGE)
%   - \CAPETCORRIGE   -> imprime le corrige          (variante corriges, INCHANGE)
%   - \CAPETMANUEL    -> imprime l'enonce ET met le corrige de cote pour l'annexe
%
% Le mode manuel stocke chaque corrige dans une macro (\csgdef), sans ecriture
% fichier ni re-tokenisation : le contenu (circuitikz, accents, %) traverse intact.
% L'annexe les rejoue ensuite dans l'ordre via \CAPETANNEXE.
\usepackage{etoolbox}

\newcounter{CAPETexo}          % numero continu d'exercice dans le manuel
\setcounter{CAPETexo}{0}
\newcounter{CAPETannexe}       % compteur de relecture en annexe

% Titre courant de la notion en cours (pose par \CAPETNOTION), sert d'intertitre annexe
\newcommand{\CAPETnotioncode}{}
\newcommand{\CAPETnotiontitre}{}

% \CAPETNOTION{code}{titre} : ouvre une notion dans le manuel (memorise le libelle,
% et marque le rang du premier exercice de la notion pour l'intertitre d'annexe).
\newcommand{\CAPETNOTION}[2]{%
  \renewcommand{\CAPETnotioncode}{#1}%
  \renewcommand{\CAPETnotiontitre}{#2}%
  \csgdef{CAPETnotionstart\the\numexpr\value{CAPETexo}+1\relax}{#1 -- #2}%
}

% ── \degree robuste ───────────────────────────────────────────────────
% Certains fichiers d'exercices utilisent \degree. Ni le preambule ni siunitx
% ne le definissent, et \providecommand local ne suffit pas : sous hyperref,
% \degree existe deja comme commande factice (\ERROR), donc \providecommand
% s'abstient et l'appel echoue a la compilation du manuel (erreur invisible
% en compilation isolee). On le (re)definit donc inconditionnellement ici,
% apres le chargement des packages, pour tout le corpus.
\AtBeginDocument{\providecommand{\degree}{}\renewcommand{\degree}{\ensuremath{^\circ}}}

% Libelle decoratif : les fichiers d'exercices utilisent aussi \CAPETSHOW dans
% leur \fancyhead et leur bloc de titre, pour afficher "(enonces)" ou "(corriges)".
% Ces appels-la ne sont PAS des exercices : en mode manuel ils ne doivent ni
% consommer un numero, ni alimenter l'annexe. On les reconnait a leur premier
% argument, strictement egal au libelle ci-dessous.
\newcommand{\CAPETLIBELLEENONCES}{(\'enonc\'es)}

\ifdefined\CAPETMANUEL
  % --- Mode manuel : enonce en place, corrige differe vers l'annexe ---
  \newcommand{\CAPETSHOW}[2]{%
    \def\CAPETargun{#1}%
    \ifx\CAPETargun\CAPETLIBELLEENONCES
      #1%            appel decoratif : rendu tel quel, hors numerotation
    \else
      \CAPETSHOWEXO{#1}{#2}%
    \fi
  }
  \newcommand{\CAPETSHOWEXO}[2]{%
    \stepcounter{CAPETexo}%
    \csgdef{CAPETcorrige\theCAPETexo}{#2}%
    \edef\CAPETtmpnum{\theCAPETexo}%
    #1%
    \par\vspace{2pt}%
    {\footnotesize\itshape\color{colExo}%
     \hyperref[capet:cor:\CAPETtmpnum]{Corrige de l'exercice \CAPETtmpnum\ $\rightarrow$ p.~\pageref{capet:cor:\CAPETtmpnum}}%
     \label{capet:eno:\CAPETtmpnum}}%
  }
  % \CAPETANNEXE : rejoue tous les corriges captures, dans l'ordre, avec intertitres
  \newcommand{\CAPETANNEXE}{%
    \setcounter{CAPETannexe}{0}%
    \loop\ifnum\value{CAPETannexe}<\value{CAPETexo}%
      \stepcounter{CAPETannexe}%
      \edef\CAPETcurnum{\theCAPETannexe}%
      \ifcsname CAPETnotionstart\CAPETcurnum\endcsname
        \par\medskip
        {\normalsize\bfseries\color{colExo}\csuse{CAPETnotionstart\CAPETcurnum}}%
        \par\nobreak\smallskip
      \fi
      \begin{tcolorbox}[exoBox, breakable,
        title={Corrige de l'exercice \CAPETcurnum\
               \hfill\normalfont\footnotesize
               \hyperref[capet:eno:\CAPETcurnum]{enonce p.~\pageref{capet:eno:\CAPETcurnum}}}]%
      \label{capet:cor:\CAPETcurnum}%
      \csuse{CAPETcorrige\CAPETcurnum}%
      \end{tcolorbox}%
    \repeat
  }
\else
  % --- Modes historiques, comportement strictement inchange ---
  \newcommand{\CAPETSHOW}[2]{\ifdefined\CAPETCORRIGE #2\else #1\fi}
  \newcommand{\CAPETANNEXE}{}
\fi

% ── Difficulté 1 à 3 (étoiles pleines, amssymb) ───────────────────────
\newcommand{\difficulte}[1]{%
  \ifcase#1\relax
  \or $\bigstar$\;niveau 1/3%
  \or $\bigstar\bigstar$\;niveau 2/3%
  \or $\bigstar\bigstar\bigstar$\;niveau 3/3%
  \fi
}

% ── Style encadré extrait de document technique (banque d'exercices) ──
\tcbset{
  docTech/.style={
    breakable,
    enhanced,
    colback=colGrayLight,
    colframe=colSeparator,
    fonttitle=\bfseries\small\itshape,
    coltitle=black,
    colbacktitle=colGrayLight,
    left=4pt, right=4pt, top=3pt, bottom=3pt,
    before skip=4pt, after skip=4pt,
  }
}

% ── Style tcolorbox de section principale ────────────────────────────
% Refonte 2026 : bandeau titre coloré arrondi + fond très pâle, sans cadre.
\tcbset{
  ficheSection/.style={
    breakable, enhanced jigsaw,
    frame hidden, boxrule=0pt, arc=1.8mm,
    colback=#1!4!white,
    colbacktitle=#1,
    coltitle=white,
    fonttitle=\sffamily\bfseries,
    left=7pt, right=7pt, top=5pt, bottom=5pt,
    toptitle=3pt, bottomtitle=3pt,
    before skip=10pt, after skip=8pt,
  }
}

% ── Style encadré question type ──────────────────────────────────────
% Carte blanche, bandeau sombre, liseré très clair.
\tcbset{
  questiontype/.style={
    breakable, enhanced jigsaw,
    colback=white,
    frame hidden, boxrule=0pt, arc=1.8mm,
    colbacktitle=colElec!78!black,
    coltitle=white,
    fonttitle=\sffamily\bfseries\small,
    left=7pt, right=7pt, top=4pt, bottom=4pt,
    toptitle=2.5pt, bottomtitle=2.5pt,
    before skip=8pt, after skip=8pt,
    borderline={0.6pt}{0pt}{colElec!12!white},
  }
}

% ── Style encadré jury ───────────────────────────────────────────────
% Bandeau ambre, fond crème, sans cadre.
\tcbset{
  jury/.style={
    breakable, enhanced jigsaw,
    colback=colJuryBg,
    frame hidden, boxrule=0pt, arc=1.8mm,
    colbacktitle=colJuryFrame,
    coltitle=white,
    fonttitle=\sffamily\bfseries,
    left=7pt, right=7pt, top=4pt, bottom=4pt,
    toptitle=3pt, bottomtitle=3pt,
    before skip=10pt, after skip=8pt,
  }
}

% ── En-tête / pied de page communs (refonte 2026) ─────────────────────
% #1 code fiche, #2 titre court, #3 couleur domaine
\newcommand{\ficheHeaders}[3]{%
  \fancyhead[L]{\small\sffamily{\bfseries\textcolor{#3}{#1}}\hspace{0.7em}\textcolor{colSlate}{#2}}%
  \fancyhead[R]{\small\sffamily\textcolor{colSlate}{CAPET SII IE\, --- \,Session 2026}}%
  \fancyfoot[C]{\small\sffamily\textcolor{colSlate}{\thepage\,/\,\pageref{LastPage}}}%
  \fancyhead[C]{}\fancyfoot[L]{}\fancyfoot[R]{}%
}

% ── Bloc en-tête de fiche (BLOC 0, refonte 2026) ──────────────────────
% #1 code  #2 titre  #3 couleur domaine  #4 domaine  #5 ligne présence  #6 sous-thèmes
\newcommand{\ficheEntete}[6]{%
  \begin{tcolorbox}[enhanced, breakable,
    frame hidden, boxrule=0pt, arc=2.5mm,
    interior style={left color=#3, right color=#3!75!black},
    left=9pt, right=9pt, top=7pt, bottom=7pt,
    before skip=0pt, after skip=6pt]
    {\color{white}\sffamily\bfseries\LARGE #1\quad #2}\\[4pt]
    {\color{#3!22!white}\sffamily\small \MakeUppercase{#4}\quad\textbar\quad #5}\\[1pt]
    {\color{#3!35!white}\sffamily\footnotesize Sous-thèmes : #6}
  \end{tcolorbox}%
}

% ── Mise en page ─────────────────────────────────────────────────────
\pagestyle{fancy}
\fancyhf{}
\renewcommand{\headrulewidth}{0.5pt}
\renewcommand{\footrulewidth}{0pt}
\AtBeginDocument{%
  \renewcommand{\headrule}{{\color{colRule}\hrule width\headwidth height 0.5pt}}%
}

\setlist[itemize]{leftmargin=1.2em, itemsep=1pt, topsep=2pt, parsep=0pt}
\setlist[enumerate]{leftmargin=1.2em, itemsep=3pt, topsep=2pt, parsep=0pt}

\setlength{\parskip}{2pt}
\setlength{\columnsep}{10pt}

\def\CAPETBEGINFILE{\begin{document}}
\def\CAPETENDFILE{\end{document}}
\else
\def\CAPETBEGINFILE{}
\def\CAPETENDFILE{}
\fi

\fancyhead[L]{\textbf{\textcolor{colExo}{T18}} - Exercices - Circuits \'electroniques, composants}
\fancyhead[R]{\small\textit{Banque d'exercices \CAPETSHOW{(\'enonc\'es)}{(corrig\'es)}}}
\fancyfoot[C]{\small\thepage\,/\,\pageref{LastPage}}

\CAPETBEGINFILE

\begin{tcolorbox}[
  enhanced, breakable,
  colback=colExo,
  colframe=colExo,
  left=6pt, right=6pt, top=5pt, bottom=5pt,
  before skip=0pt, after skip=4pt
]
  {\color{white}\bfseries\large T18 \quad Circuits \'electroniques, composants (diode, transistor)}\\[2pt]
  {\color{white}\small Banque d'exercices -- 6 exercices \CAPETSHOW{(\'enonc\'es)}{(corrig\'es)}}
\end{tcolorbox}

\begin{tcolorbox}[exoBox,
  title={Exercice 1 -- \difficulte{1} \hfill\normalfont\footnotesize 2022, \'epreuve \'ecrite disciplinaire, questions 26 et 27}
]
\CAPETSHOW{
Le module de surveillance TRT200 est secouru par des accumulateurs. Le courant de charge/d\'echarge de ces accumulateurs traverse une r\'esistance $R_0$ de faible valeur, dont la tension est appliqu\'ee \`a l'entr\'ee analogique ADC1 du convertisseur analogique-num\'erique (CAN) du microcontr\^oleur. Une diode $D_1$ relie cette entr\'ee au $\SI{0}{\volt}$, anode c\^ot\'e $\SI{0}{\volt}$, cathode c\^ot\'e entr\'ee du CAN. La tension de seuil des diodes est prise \'egale \`a \SI{0,6}{\volt} et la tension de r\'ef\'erence du CAN vaut \SI{3,5}{\volt}.

\begin{center}
\begin{circuitikz}[european, scale=0.8]
\draw (0,3) node[left]{accumulateurs} -- (1.2,3)
to[R, l=$R_0$, i>^=$i_{ch}$] (4.2,3) node[circ]{}
-- (5.6,3) node[ocirc]{} node[above=5pt]{\footnotesize vers ADC1};
\draw (4.2,3) to[empty diode, l_=$D_1$, invert] (4.2,0.6) -- (4.2,0.4) node[ground]{};
\draw (1.2,3) -- (1.2,0.6) -- (1.2,0.4) node[ground]{};
\draw[-latex] (5.6,2.7) -- (5.6,0.9) node[midway, right]{\footnotesize $v_{ADC1}$};
\end{circuitikz}
\end{center}

\begin{enumerate}[leftmargin=1.5em]
  \item Quelle fonction $R_0$ remplit-elle dans la mesure de la charge des accumulateurs~? Pr\'eciser la grandeur physique r\'eellement image de la tension pr\'elev\'ee \`a ses bornes.
  \item Les accumulateurs sont d'abord en \emph{charge}. Indiquer le sens du courant $i_{ch}$ dans $R_0$, en d\'eduire le signe de $v_{ADC1}$, puis l'\'etat (passante / bloqu\'ee) de $D_1$.
  \item Les accumulateurs d\'ebitent ensuite vers le syst\`eme (\emph{d\'echarge}). Reprendre le raisonnement de la question pr\'ec\'edente et en d\'eduire la valeur limite que peut atteindre $v_{ADC1}$.
  \item Contre quel type d'agression \'electrique $D_1$ prot\`ege-t-elle donc l'entr\'ee du CAN~? Justifier en une phrase pourquoi cette agression serait critique pour le microcontr\^oleur.
\end{enumerate}

}{
\textbf{Corrig\'e.}

\begin{enumerate}[leftmargin=1.5em]
  \item $R_0$ est une \textbf{r\'esistance de shunt} : elle est ins\'er\'ee en s\'erie sur le circuit de charge/d\'echarge et la tension pr\'elev\'ee \`a ses bornes, $v = R_0\,i_{ch}$, est une \textbf{image du courant} de charge des accumulateurs (et non de leur tension, surveill\'ee s\'epar\'ement sur ADC0). La mesure d'un courant est ainsi ramen\'ee \`a une mesure de tension, seule grandeur que le CAN sait num\'eriser.
  \item En charge, le courant circule de la borne $+$ vers la borne $-$ des accumulateurs, c'est-\`a-dire dans le sens positif de $i_{ch}$ sur le sch\'ema : $v_{ADC1} = R_0\,i_{ch} > 0$. La diode $D_1$, dont l'anode est au $\SI{0}{\volt}$ et la cathode sur l'entr\'ee, voit alors $V_{AK} = 0 - v_{ADC1} < 0$ : elle est \textbf{bloqu\'ee} et n'a aucune influence sur la mesure.
  \item En d\'echarge, le courant change de sens : $i_{ch} < 0$ et $v_{ADC1} = R_0\,i_{ch}$ devient \textbf{n\'egative}. D\`es que $v_{ADC1}$ descend sous $-V_{seuil}$, on a $V_{AK} = 0 - v_{ADC1} > \SI{0,6}{\volt}$ : $D_1$ devient \textbf{passante} et \'ecr\^ete la tension. L'entr\'ee du CAN ne peut donc pas descendre en dessous de
  \[ v_{ADC1,min} = -V_{seuil} = \SI{-0,6}{\volt} \]
  \item $D_1$ prot\`ege l'entr\'ee analogique contre une \textbf{tension n\'egative}, apparue par simple inversion du sens du courant lors de la d\'echarge. C'est critique car la plage d'entr\'ee admissible du CAN est born\'ee par ses rails d'alimentation ($\SI{0}{\volt}$ \`a $V_{ref} = \SI{3,5}{\volt}$) : une tension n\'egative significative polarise en direct les diodes de structure internes du microcontr\^oleur, injecte un courant dans le substrat et peut d\'et\'eriorer d\'efinitivement l'\'etage d'entr\'ee.
\end{enumerate}

\textbf{Conclusion.} L'association $R_0 + D_1$ satisfait deux exigences \`a la fois : $R_0$ rend le courant d'accumulateur observable par un CAN qui ne mesure que des tensions, et $D_1$ garantit que la grandeur pr\'esent\'ee \`a ADC1 ne franchit jamais $\SI{-0,6}{\volt}$ malgr\'e l'inversion de sens du courant en d\'echarge -- donc reste dans la plage admissible de l'entr\'ee analogique du PIC. Sans cette diode, le simple passage du TRT200 sur son alimentation de secours suffirait \`a d\'et\'eriorer le microcontr\^oleur de surveillance.
}
\end{tcolorbox}

\begin{tcolorbox}[exoBox,
  title={Exercice 2 -- \difficulte{1} \hfill\normalfont\footnotesize 2013, \'epreuve commune de synth\`ese de sciences industrielles, question 48}
]
\CAPETSHOW{
Une installation photovolta\"ique raccord\'ee au r\'eseau comporte 24 onduleurs de r\'ef\'erence EI3300. Chaque onduleur transforme la tension continue issue des panneaux en une tension alternative sinuso\"idale de valeur efficace \SI{230}{\volt} et de fr\'equence \SI{50}{\hertz}, compatible avec le r\'eseau EDF. Ses interrupteurs de puissance sont des IGBT.

\begin{center}
\begin{tikzpicture}[scale=0.9, >=Stealth, node distance=1cm]
  \node[draw, minimum width=2cm, minimum height=1.2cm] (pv) {champ PV};
  \node[draw, minimum width=2.6cm, minimum height=1.2cm, right=1.5cm of pv] (ond) {onduleur (IGBT)};
  \node[draw, minimum width=2cm, minimum height=1.2cm, right=2.6cm of ond] (res) {r\'eseau EDF};
  \draw[-latex] (pv) -- node[above]{$U_{DC}$} (ond);
  \draw[-latex] (ond) -- node[above, font=\footnotesize]{\SI{230}{\volt}, \SI{50}{\hertz}} (res);
\end{tikzpicture}
\end{center}

\begin{enumerate}[leftmargin=1.5em]
  \item Pr\'eciser le r\^ole fonctionnel de l'onduleur dans cette cha\^ine (deux fonctions \`a distinguer).
  \item Quel type de commande est utilis\'e pour les transistors de puissance~? Pourquoi ce choix est-il pr\'ef\'er\'e \`a une commande dite \guillemotleft\ pleine onde \guillemotright\ pour une injection sur le r\'eseau~?
  \item La tension continue en entr\'ee de l'onduleur (c\^ot\'e panneaux) est sup\'erieure \`a l'amplitude cr\^ete de la tension r\'eseau. Justifier pourquoi des transistors command\'es (IGBT), et non de simples diodes, sont indispensables ici.
\end{enumerate}

}{
\textbf{Corrig\'e.}

\begin{enumerate}[leftmargin=1.5em]
  \item L'onduleur assure deux fonctions : convertir la tension continue issue des panneaux en tension alternative, et rendre cette tension compatible avec le r\'eseau (valeur efficace \SI{230}{\volt}, fr\'equence \SI{50}{\hertz}, synchronisation en phase) afin de permettre l'injection de l'\'energie produite chez EDF.
  \item La commande est de type \textbf{MLI} (modulation de largeur d'impulsion, ou PWM) : le d\'ecoupage \`a haute fr\'equence permet d'obtenir en sortie une tension dont le fondamental est quasi-sinuso\"idal, avec un spectre d'harmoniques repouss\'e vers les hautes fr\'equences (donc facile \`a filtrer), contrairement \`a une commande pleine onde dont les harmoniques de rang faible restent proches du fondamental. C'est n\'ecessaire pour respecter les exigences de qualit\'e de l'onde inject\'ee sur le r\'eseau.
  \item Un pont \`a diodes est un convertisseur unidirectionnel : le courant ne peut circuler que dans un seul sens impos\'e par les diodes, ce qui interdit tout transfert command\'e de puissance active du c\^ot\'e continu vers le c\^ot\'e alternatif. Des transistors command\'es (IGBT) permettent au contraire de piloter le sens et l'instant du transfert d'\'energie : c'est cette r\'eversibilit\'e command\'ee qui rend possible l'injection de puissance sur le r\'eseau.
\end{enumerate}

\textbf{Conclusion.} Structure \`a interrupteurs command\'es et commande MLI ne sont pas ici des choix de confort~: ce sont les deux conditions pour que les 24 onduleurs EI3300 satisfassent le crit\`ere de raccordement au r\'eseau de distribution (onde \SI{230}{\volt} / \SI{50}{\hertz} synchronis\'ee et faiblement distordue), sans lequel l'\'energie produite ne pourrait pas \^etre revendue.
}
\end{tcolorbox}

\begin{tcolorbox}[exoBox,
  title={Exercice 3 -- \difficulte{2} \hfill\normalfont\footnotesize exercice in\'edit}
]
\CAPETSHOW{
Une diode au silicium est mont\'ee en s\'erie avec une r\'esistance $R = \SI{220}{\ohm}$, l'ensemble \'etant aliment\'e par une source id\'eale $E = \SI{5}{\volt}$.

\begin{center}
\begin{circuitikz}[european, scale=0.85]
\draw
(0,0) to[V, l=$E$] (0,3)
to[R, l=$R$] (3,3)
to[empty diode, l=$D$, *-*] (3,0)
-- (0,0);
\end{circuitikz}
\end{center}

La caract\'eristique statique $i(V_{AK})$ de la diode est donn\'ee ci-dessous (axes gradu\'es en volts et milliamp\`eres).

\begin{center}
\begin{tikzpicture}[scale=1, >=Stealth]
  \draw[-latex] (0,0) -- (6,0) node[right]{$V_{AK}$ (V)};
  \draw[-latex] (0,0) -- (0,5) node[above]{$i$ (mA)};
  \foreach \x in {1,2,3,4,5} \draw (\x,0.05) -- (\x,-0.05) node[below]{\x};
  \foreach \y/\lbl in {1/5,2/10,3/15,4/20} \draw (0.05,\y) -- (-0.05,\y) node[left]{\lbl};
  \draw[thick, colExo] (0,0) -- (0.54,0) .. controls (0.60,0) and (0.64,3.3) .. (0.68,4.55) -- (0.82,4.85);
  \node[colExo, right] at (1.05,4.5) {caract\'eristique diode};
\end{tikzpicture}
\end{center}

\begin{enumerate}[leftmargin=1.5em]
  \item \'Etablir l'\'equation de la droite de charge $i = f(V_{AK})$ impos\'ee par le circuit ext\'erieur ($E$ et $R$).
  \item Tracer cette droite de charge sur le graphe et d\'eterminer graphiquement le point de fonctionnement $Q$ (coordonn\'ees approximatives $V_{AK,Q}$ et $I_Q$).
  \item Retrouver $I_Q$ par le calcul en utilisant le mod\`ele \`a seuil $V_{seuil} = \SI{0,6}{\volt}$, et comparer au r\'esultat graphique.
\end{enumerate}

}{
\textbf{Corrig\'e.}

\begin{enumerate}[leftmargin=1.5em]
  \item Loi des mailles : $E = R\cdot i + V_{AK}$, soit :
  \[ i = \dfrac{E - V_{AK}}{R} \]
  C'est l'\'equation d'une droite d'ordonn\'ee \`a l'origine $E/R = 5/220 \approx \SI{22,7}{\milli\ampere}$ (en $V_{AK}=0$) et d'abscisse \`a l'origine $E = \SI{5}{\volt}$ (en $i=0$), de pente $-1/R$.

  \item On trace la droite de charge en joignant ces deux points ; le point de fonctionnement $Q$ est son intersection avec la caract\'eristique statique de la diode :

  \begin{center}
  \begin{tikzpicture}[scale=1, >=Stealth]
    \draw[-latex] (0,0) -- (6,0) node[right]{$V_{AK}$ (V)};
    \draw[-latex] (0,0) -- (0,5) node[above]{$i$ (mA)};
    \foreach \x in {1,2,3,4,5} \draw (\x,0.05) -- (\x,-0.05) node[below]{\x};
    \foreach \y/\lbl in {1/5,2/10,3/15,4/20} \draw (0.05,\y) -- (-0.05,\y) node[left]{\lbl};
    \draw[thick, colExo] (0,0) -- (0.54,0) .. controls (0.60,0) and (0.64,3.3) .. (0.68,4.55) -- (0.82,4.85);
    \draw[thick, dashed] (0,4.55) -- (5,0) node[pos=0.55, above]{droite de charge};
    \filldraw (0.664,3.947) circle (1.7pt) node[right=3pt, yshift=3pt]{$Q$};
    \draw[dotted] (0.664,0) -- (0.664,3.947);
    \draw[dotted] (0,3.947) -- (0.664,3.947);
  \end{tikzpicture}
  \end{center}
  Graphiquement, l'intersection $Q$ de la droite de charge et de la caract\'eristique se situe vers $V_{AK,Q} \approx \SI{0,66}{\volt}$, $I_Q \approx \SI{20}{\milli\ampere}$.

  \item Avec le mod\`ele \`a seuil, $V_{AK} = V_{seuil} = \SI{0,6}{\volt}$ quel que soit $i>0$, donc :
  \[ I_Q = \dfrac{E - V_{seuil}}{R} = \dfrac{5 - 0{,}6}{220} \approx \SI{20,0}{\milli\ampere} \]
  Ce r\'esultat est tr\`es proche de la lecture graphique : le mod\`ele \`a seuil est une bonne approximation d\`es lors que $R$ est suffisamment grande devant la r\'esistance dynamique de la diode au voisinage de $Q$ (la droite de charge est alors peu pentue, et $V_{AK,Q}$ reste proche de $V_{seuil}$).
\end{enumerate}

\textbf{Conclusion.} L'\'ecart entre les deux m\'ethodes reste inf\'erieur au pour-cent sur $I_Q$ : pour dimensionner le circuit (choix du calibre de $R$, v\'erification du courant direct admissible de la diode), le mod\`ele \`a seuil suffit et dispense du trac\'e graphique ; la r\'esolution graphique ne redevient n\'ecessaire que si la r\'esistance dynamique de la diode cesse d'\^etre n\'egligeable devant $R$.
}
\end{tcolorbox}

\begin{tcolorbox}[exoBox,
  title={Exercice 4 -- \difficulte{2} \hfill\normalfont\footnotesize exercice in\'edit, inspir\'e du convertisseur CV3-1 (2025, \'epreuve \'ecrite disciplinaire)}
]
\CAPETSHOW{
On \'etudie un hacheur inspir\'e du convertisseur CV3-1 d'un sous-r\'eseau d'\'electrification aliment\'e par pile \`a combustible : il recharge ici une batterie \`a partir du bus continu $V_{DC}$. Le sch\'ema simplifi\'e comporte un transistor $Q_1$ command\'e (rapport cyclique $\alpha$, p\'eriode $T$) et une diode $D_1$, tous deux suppos\'es parfaits, une inductance $L_1$ parcourue par $i_{L1}$, et la batterie mod\'elis\'ee par une f.\'e.m. $E_b$. Les tensions $V_{DC}$, $u_{L1}$, $v_{Q1}$, $v_{D1}$ et $E_b$, ainsi que les courants $i_{L1}$, $i_{Q1}$ et $i_{D1}$, sont orient\'es selon les fl\`eches du sch\'ema ($u_{L1}$, $v_{Q1}$ et $v_{D1}$ en convention r\'ecepteur par rapport aux courants correspondants).

\begin{center}
\begin{circuitikz}[european, scale=0.8]
\draw
(0,0) to[V] (0,4)
(0,4) to[normal open switch, l=$Q_1$] (3,4)
(3,0) to[empty diode, l=$D_1$] (3,4)
(3,0) -- (0,0)
(3,4) to[cute inductor, l=$L_1$] (6,4)
(6,4) to[battery2] (6,0)
(6,0) -- (3,0)
;
\draw (3,4) node[circ]{};
\draw (3,0) node[circ]{};
\draw[-latex] (-0.95,0.8) -- (-0.95,3.2) node[midway, left]{$V_{DC}$};
\draw[-latex] (6.95,0.8) -- (6.95,3.2) node[midway, right]{$E_b$};
\draw[-latex, dashed] (1.5,4.75) -- (1.5,4.15);
\node[above] at (1.5,4.75) {\footnotesize commande ($\alpha$, $T$)};
\draw[-latex] (5.7,6.5) -- (3.3,6.5) node[midway, above]{$u_{L1}$};
\draw[-latex] (4.3,3.55) -- (5.3,3.55) node[midway, below]{$i_{L1}$};
\draw[-latex] (2.6,6.5) -- (0.4,6.5) node[midway, above]{\footnotesize $v_{Q1}$};
\draw[-latex] (0.35,3.6) -- (1.05,3.6) node[midway, below]{\footnotesize $i_{Q1}$};
\draw[-latex] (1.75,1.15) -- (1.75,2.85) node[midway, left]{\footnotesize $v_{D1}$};
\draw[-latex] (3.75,1.35) -- (3.75,2.65) node[midway, right]{\footnotesize $i_{D1}$};
\end{circuitikz}
\end{center}

\begin{enumerate}[leftmargin=1.5em]
  \item Pour chacune des quatre combinaisons possibles des \'etats de $Q_1$ (ouvert/ferm\'e) et $D_1$ (bloqu\'ee/passante), indiquer si elle est physiquement viable, en justifiant \`a partir du sens de conduction impos\'e \`a $D_1$ et de la continuit\'e du courant $i_{L1}$ dans l'inductance.
  \item \'Ecrire les lois g\'en\'erales de Kirchhoff (loi des n\oe uds et loi des mailles) qui lient entre elles les grandeurs fl\'ech\'ees $V_{DC}$, $u_{L1}$, $v_{Q1}$, $v_{D1}$, $E_b$ d'une part, et les courants $i_{L1}$, $i_{Q1}$, $i_{D1}$ d'autre part. Ces relations doivent \^etre valables \`a tout instant, ind\'ependamment de l'\'etat des interrupteurs.
  \item Justifier que la valeur moyenne de la tension aux bornes de $L_1$ est nulle en r\'egime \'etabli, puis en d\'eduire la relation entre $E_b$, $V_{DC}$ et $\alpha$.
\end{enumerate}

}{
\textbf{Corrig\'e.}

\begin{enumerate}[leftmargin=1.5em]
  \item Les quatre combinaisons :
  \begin{itemize}
    \item $Q_1$ ferm\'e, $D_1$ bloqu\'ee : \textbf{viable}. $Q_1$ ferm\'e impose le n\oe ud commun au potentiel $V_{DC}$ ; $D_1$, dont la cathode est reli\'ee \`a ce n\oe ud et l'anode au rail $\SI{0}{\volt}$, voit alors une tension inverse et reste naturellement bloqu\'ee.
    \item $Q_1$ ferm\'e, $D_1$ passante : \textbf{impossible}. $D_1$ passante imposerait le n\oe ud \`a un potentiel proche de $0$, en contradiction directe avec $Q_1$ ferm\'e qui l'impose \`a $V_{DC}$ -- cela court-circuiterait la source.
    \item $Q_1$ ouvert, $D_1$ bloqu\'ee : \textbf{impossible} en conduction continue. Le courant $i_{L1}$ dans une inductance ne peut subir de discontinuit\'e ; si $Q_1$ s'ouvre et qu'aucun chemin n'est disponible pour $i_{L1}$, une surtension destructrice appara\^itrait aux bornes de $L_1$ ($di/dt$ infini). $D_1$ doit donc devenir passante pour assurer la continuit\'e du courant.
    \item $Q_1$ ouvert, $D_1$ passante : \textbf{viable}. C'est la phase de roue libre : $D_1$ assure la continuit\'e de $i_{L1}$ pendant le blocage de $Q_1$.
  \end{itemize}
  \item \textbf{Loi des n\oe uds} au point de commutation (n\oe ud commun \`a $Q_1$, $D_1$ et $L_1$) : le courant apport\'e par $Q_1$ et celui apport\'e par $D_1$ se rejoignent dans l'inductance, donc \`a tout instant
  \[ i_{L1} = i_{Q1} + i_{D1} \]
  \textbf{Loi des mailles.} Maille source-transistor-diode (source $V_{DC}$, $Q_1$, $D_1$, retour par le rail $\SI{0}{\volt}$) :
  \[ V_{DC} = v_{Q1} + v_{D1} \]
  Maille diode-inductance-batterie ($D_1$, $L_1$, $E_b$) :
  \[ v_{D1} = u_{L1} + E_b \quad\text{avec}\quad u_{L1} = L_1\dfrac{di_{L1}}{dt} \]
  Par \'elimination de $v_{D1}$, on obtient la relation valable \`a tout instant :
  \[ V_{DC} - v_{Q1} = L_1\dfrac{di_{L1}}{dt} + E_b \]
  Ces trois relations sont bien ind\'ependantes de l'\'etat des interrupteurs : ce sont les \'etats de $Q_1$ et $D_1$ qui viennent ensuite fixer les valeurs particuli\`eres de $v_{Q1}$ et $v_{D1}$ (interrupteur ferm\'e $\Rightarrow$ tension nulle, interrupteur ouvert/bloqu\'e $\Rightarrow$ courant nul).
  \item En r\'egime p\'eriodique \'etabli, $i_{L1}(T) = i_{L1}(0)$, donc :
  \[ \left\langle L_1\dfrac{di_{L1}}{dt}\right\rangle = \dfrac{L_1}{T}\big[i_{L1}(T)-i_{L1}(0)\big] = 0 \]
  On moyenne alors sur une p\'eriode la relation \'etablie \`a la question pr\'ec\'edente, $V_{DC} - v_{Q1} = L_1\,di_{L1}/dt + E_b$. D'apr\`es la question 1, seuls deux \'etats sont visit\'es : $Q_1$ ferm\'e ($v_{Q1}=0$) pendant $\alpha T$, puis $Q_1$ ouvert avec $D_1$ passante ($v_{D1}=0$, donc $v_{Q1}=V_{DC}$) pendant $(1-\alpha)T$. La grandeur $V_{DC}-v_{Q1}$ vaut donc $V_{DC}$ pendant $\alpha T$ et $0$ pendant $(1-\alpha)T$, soit une valeur moyenne $\alpha V_{DC}$ :
  \[ \alpha V_{DC} = E_b + \left\langle L_1\dfrac{di_{L1}}{dt}\right\rangle = E_b \quad\Longrightarrow\quad E_b = \alpha \cdot V_{DC} \]
\end{enumerate}

\textbf{Conclusion.} Comme $0 \leq \alpha \leq 1$, la structure impose $E_b \leq V_{DC}$ : il s'agit d'un hacheur d\'evolteur (Buck), dont le rapport cyclique est le seul param\`etre de r\'eglage de la tension de charge de la batterie. Le crit\`ere de dimensionnement en d\'ecoule directement : le bus continu doit rester au-dessus de la tension de fin de charge de la batterie, faute de quoi aucune valeur de $\alpha$ ne permet d'atteindre la consigne.

\emph{Remarque -- lien avec le convertisseur CV3-1 du sujet 2025.} Dans le convertisseur r\'eel, les r\^oles de $Q_1$ et $D_1$ sont \'echang\'es ($Q_1$ en parall\`ele sur le rail, $D_1$ en s\'erie vers le bus) : la m\^eme m\'ethode ($\langle u_{L1}\rangle = 0$) conduit alors \`a la relation \'el\'evatrice $V_{DC} = E_b/(1-\alpha)$, coh\'erente avec le point de repos donn\'e par le sujet, $E_{b0}/V_{DC0} = 1-\alpha_0 = 48/230$.
}
\end{tcolorbox}

\begin{tcolorbox}[exoBox,
  title={Exercice 5 -- \difficulte{2} \hfill\normalfont\footnotesize lecture de document technique -- exercice in\'edit}
]
\CAPETSHOW{
Un microcontr\^oleur commande, via sa sortie logique $V_{cmd}$ (0 ou $\SI{5}{\volt}$), la bobine d'un relais \'electrom\'ecanique (r\'esistance $R_{relais} = \SI{100}{\ohm}$, aliment\'ee sous $V_{cc} = \SI{12}{\volt}$) par l'interm\'ediaire d'un transistor bipolaire NPN mont\'e en \'emetteur commun. Une r\'esistance de base $R_B$ limite le courant de base $I_B$, et une diode de roue libre $D$ est plac\'ee en antiparall\`ele sur la bobine. Le composant retenu est un \textbf{BC337-25} en bo\^itier TO-92, dont un extrait de la fiche constructeur est reproduit ci-dessous.

\begin{center}
\begin{circuitikz}[european, scale=0.85]
\draw (0,0) node[npn] (Q1) {};
\node[right=8pt] at (Q1.E) {$Q_1$};
\draw (Q1.B) -- ++(-2,0) to[R, l=$R_B$] ++(-2,0) node[left]{$V_{cmd}$};
\draw (Q1.E) -- ++(0,-1) node[ground]{};
\draw (Q1.C) -- ++(0,1.5) coordinate (A);
\draw (A) node[circ]{} to[R, l_=$R_{relais}$] ++(3,0) coordinate (T1);
\draw (A) -- ++(0,1.8) coordinate (A2);
\draw (A2) to[empty diode, l=$D$] ++(3,0) coordinate (T2);
\draw (T1) node[circ]{} -- (T2);
\draw (T1) to[V, l_=$V_{cc}$] ++(0,-2.9) node[ground]{};
\end{circuitikz}
\end{center}

\begin{tcolorbox}[docTech, title={Document technique -- BC337-25, NPN General Purpose Amplifier (extrait)}]
\footnotesize
\textbf{Absolute Maximum Ratings} ($T_A = \SI{25}{\degreeCelsius}$)

\begin{center}
\begin{tabular}{llcc}
\toprule
\textbf{Symbol} & \textbf{Parameter} & \textbf{Value} & \textbf{Units} \\
\midrule
$V_{CEO}$ & Collector-Emitter Voltage & 45 & V \\
$V_{CES}$ & Collector-Base Voltage & 50 & V \\
$V_{EBO}$ & Emitter-Base Voltage & 5,0 & V \\
$I_C$ & Collector Current -- Continuous & 1,0 & A \\
$P_D$ & Total Device Dissipation ($T_A = \SI{25}{\degreeCelsius}$) & 625 & mW \\
$T_J$, $T_{stg}$ & Operating and Storage Junction Temp. Range & $-55$ \`a $+150$ & \si{\degreeCelsius} \\
\bottomrule
\end{tabular}
\end{center}

\textbf{Electrical Characteristics} -- \emph{ON characteristics} ($T_A = \SI{25}{\degreeCelsius}$)

\begin{center}
\begin{tabular}{llccc}
\toprule
\textbf{Symbol} & \textbf{Parameter} / \textbf{Test Conditions} & \textbf{Min} & \textbf{Max} & \textbf{Units} \\
\midrule
$h_{FE}$ & DC Current Gain -- $V_{CE} = \SI{1,0}{\volt}$, $I_C = \SI{100}{\milli\ampere}$ & 160 & 400 & \\
         & \hspace{1.2em} $V_{CE} = \SI{1,0}{\volt}$, $I_C = \SI{500}{\milli\ampere}$ & 40 & -- & \\
$V_{CE(sat)}$ & Collector-Emitter Saturation Voltage & -- & 0,7 & V \\
              & \hspace{1.2em} $I_C = \SI{500}{\milli\ampere}$, $I_B = \SI{50}{\milli\ampere}$ & & & \\
$V_{BE(on)}$ & Base-Emitter On Voltage -- $V_{CE} = \SI{1,0}{\volt}$, $I_C = \SI{500}{\milli\ampere}$ & -- & 1,2 & V \\
\bottomrule
\end{tabular}
\end{center}
\end{tcolorbox}

\begin{enumerate}[leftmargin=1.5em]
  \item \textbf{Relever dans l'extrait} les trois valeurs n\'ecessaires au dimensionnement de cette commande : la tension collecteur-\'emetteur maximale admissible, le courant collecteur continu maximal et la valeur \emph{minimale} garantie du gain statique $h_{FE}$ pour $I_C = \SI{100}{\milli\ampere}$.
  \item Calculer le courant de collecteur $I_{C,sat}$ appel\'e par la bobine lorsque le transistor est satur\'e, en prenant $V_{CE(sat)}$ \'egal \`a la valeur maximale annonc\'ee par le constructeur. V\'erifier que $I_C$ et $V_{CE}$ restent dans les limites relev\'ees \`a la question 1.
  \item On impose un facteur de sursaturation $k = \beta_{min}\,I_B / I_{C,sat} = 2$. En prenant pour $\beta_{min}$ la valeur relev\'ee \`a la question 1 et $V_{BE} = \SI{0,6}{\volt}$, calculer le courant de base n\'ecessaire puis la valeur maximale de $R_B$ qui garantit cette sursaturation. Proposer une valeur normalis\'ee de la s\'erie E12.
  \item V\'erifier que la puissance dissip\'ee par le transistor en conduction reste inf\'erieure \`a la dissipation maximale $P_D$ du bo\^itier.
  \item Sans la diode de roue libre $D$, quelle grandeur du document technique serait d\'epass\'ee \`a l'ouverture de $Q_1$~? Conclure sur le caract\`ere obligatoire de $D$.
\end{enumerate}

}{
\textbf{Corrig\'e.}

\begin{enumerate}[leftmargin=1.5em]
  \item Lecture directe de l'extrait : $V_{CEO} = \SI{45}{\volt}$ (tension collecteur-\'emetteur maximale), $I_{C,max} = \SI{1,0}{\ampere}$ en continu, et $h_{FE} \geq 160$ pour $V_{CE} = \SI{1,0}{\volt}$, $I_C = \SI{100}{\milli\ampere}$ (classe $-25$). On retiendra donc $\beta_{min} = 160$.
  \item En saturation, la maille collecteur donne $V_{cc} = R_{relais}\,I_{C,sat} + V_{CE(sat)}$ :
  \[ I_{C,sat} = \dfrac{V_{cc} - V_{CE(sat)}}{R_{relais}} = \dfrac{12 - 0{,}7}{100} = \SI{113}{\milli\ampere} \]
  Ce courant est tr\`es inf\'erieur au calibre continu $I_{C,max} = \SI{1,0}{\ampere}$ (marge $\times 8{,}8$). \`A l'\'etat bloqu\'e, le transistor supporte au plus $V_{CE} = V_{cc} = \SI{12}{\volt}$, tr\`es inf\'erieur \`a $V_{CEO} = \SI{45}{\volt}$. Les deux limites du document technique sont donc respect\'ees.

  \emph{Remarque}~: le point de fonctionnement $I_C \approx \SI{113}{\milli\ampere}$ est proche de la condition de test du gain ($\SI{100}{\milli\ampere}$), ce qui l\'egitime l'emploi de $\beta_{min} = 160$.
  \item La condition de sursaturation impose $\beta_{min}\,I_B = k\,I_{C,sat}$, soit :
  \[ I_B = \dfrac{k\,I_{C,sat}}{\beta_{min}} = \dfrac{2\times 113\times 10^{-3}}{160} \approx \SI{1,41}{\milli\ampere} \]
  La maille de base donne $V_{cmd} = R_B\,I_B + V_{BE}$, d'o\`u :
  \[ R_{B,max} = \dfrac{V_{cmd} - V_{BE}}{I_B} = \dfrac{5 - 0{,}6}{1{,}41\times 10^{-3}} \approx \SI{3,12}{\kilo\ohm} \]
  Toute r\'esistance inf\'erieure convient (elle augmente $I_B$, donc $k$). On retient la valeur normalis\'ee E12 imm\'ediatement inf\'erieure : $R_B = \SI{2,7}{\kilo\ohm}$, qui donne $I_B = (5-0{,}6)/2\,700 \approx \SI{1,63}{\milli\ampere}$, soit $k = 160\times 1{,}63/113 \approx 2{,}3$.
  \item En conduction, la puissance dissip\'ee dans le transistor vaut essentiellement celle du jonction collecteur-\'emetteur :
  \[ P = V_{CE(sat)}\cdot I_{C,sat} + V_{BE}\cdot I_B = 0{,}7\times 0{,}113 + 0{,}6\times 1{,}63\times 10^{-3} \approx \SI{80}{\milli\watt} \]
  soit environ \SI{13}{\percent} de la dissipation maximale $P_D = \SI{625}{\milli\watt}$ du bo\^itier TO-92 : aucun radiateur n'est n\'ecessaire.
  \item \`A l'ouverture de $Q_1$, la bobine du relais impose la continuit\'e de son courant : la tension $v = L\,di/dt$ qui appara\^it \`a ses bornes s'ajoute \`a $V_{cc}$ sur le collecteur et d\'epasse tr\`es largement la limite $V_{CEO} = \SI{45}{\volt}$ relev\'ee dans le document technique -- c'est un claquage par avalanche du transistor. La diode $D$, en offrant un chemin de roue libre, borne cette tension \`a $V_{cc} + V_{seuil} \approx \SI{12,6}{\volt}$.
\end{enumerate}

\textbf{Conclusion.} Avec $R_B = \SI{2,7}{\kilo\ohm}$, les quatre limites du constructeur sont simultan\'ement respect\'ees~: $I_C = \SI{113}{\milli\ampere} \ll \SI{1}{\ampere}$, $V_{CE} \leq \SI{12,6}{\volt} \ll \SI{45}{\volt}$ gr\^ace \`a $D$, $P = \SI{80}{\milli\watt} \ll \SI{625}{\milli\watt}$, et le facteur de sursaturation $k \approx 2{,}3$ garantit la saturation m\^eme pour un exemplaire de gain minimal. Le BC337-25 est donc valid\'e pour cette commande de relais, \`a la condition imp\'erative de conserver la diode de roue libre, seul \'el\'ement qui emp\^eche le d\'epassement de $V_{CEO}$.
}
\end{tcolorbox}

\begin{tcolorbox}[exoBox,
  title={Exercice 6 -- \difficulte{3} \hfill\normalfont\footnotesize exercice in\'edit}
]
\CAPETSHOW{
Un redresseur monophas\'e demi-onde \`a un seul thyristor (montage P1) alimente une charge r\'esistive pure $R$ \`a partir d'une tension sinuso\"idale $v(t) = V_m\sin(\theta)$, avec $\theta = \omega t$. Le thyristor $Th$ est amorc\'e avec un retard angulaire $\alpha$ \`a chaque alternance positive, et se bloque naturellement lorsque le courant (donc la tension, charge r\'esistive) s'annule.

\begin{center}
\begin{tikzpicture}[scale=0.85, >=Stealth]
  \draw[-latex] (0,-2.1) -- (0,2.1) node[above]{$u_S,\,v$};
  \draw[-latex] (-0.3,0) -- (12.3,0) node[right]{$\theta$};
  \draw[dashed, gray] plot[domain=0:465, samples=140] ({\x/40}, {1.5*sin(\x)});
  \draw[thick, colExo] plot[domain=60:180, samples=60] ({\x/40}, {1.5*sin(\x)});
  \draw[thick, colExo] (4.5,0) -- (10.5,0);
  \draw[thick, colExo] plot[domain=420:465, samples=30] ({\x/40}, {1.5*sin(\x)});
  \draw[dotted] (1.5,0) -- (1.5,{1.5*sin(60)}) node[above]{} ; \node[below] at (1.5,-0.1) {$\alpha$};
  \draw[dotted] (4.5,0) -- (4.5,1.6); \node[below] at (4.5,-0.1) {$\pi$};
  \draw[dotted] (9,0) -- (9,1.6); \node[below] at (9,-0.1) {$2\pi$};
  \draw[dotted] (10.5,0) -- (10.5,{1.5*sin(60)}); \node[below] at (10.5,-0.1) {$2\pi{+}\alpha$};
  \node[left] at (0,1.5) {$V_m$};
\end{tikzpicture}
\end{center}

\begin{enumerate}[leftmargin=1.5em]
  \item D'apr\`es le chronogramme, pr\'eciser l'intervalle angulaire de conduction de $Th$ sur une p\'eriode $[0\,;\,2\pi]$, et justifier pourquoi le blocage a lieu pr\'ecis\'ement en $\theta = \pi$ pour une charge r\'esistive.
  \item Montrer que la valeur moyenne de $u_S(\theta)$ s'\'ecrit $\displaystyle\langle u_S\rangle = \dfrac{V_m}{2\pi}(1+\cos\alpha)$.
  \item Calculer $\langle u_S\rangle$ pour $V_m = \SI{325}{\volt}$ (secteur \SI{230}{\volt} efficace) et $\alpha = \SI{60}{\degree}$. Comparer au cas $\alpha = 0$ et conclure sur le r\^ole de l'angle d'amor\c cage.
\end{enumerate}

}{
\textbf{Corrig\'e.}

\begin{enumerate}[leftmargin=1.5em]
  \item $Th$ conduit sur $\theta\in[\alpha\,;\,\pi]$. Pour $\theta > \pi$, $v(\theta) = V_m\sin\theta$ devient n\'egative ; avec une charge r\'esistive pure, le courant $i(\theta) = u_S(\theta)/R$ est proportionnel \`a la tension et s'annulerait donc en changeant de signe en $\theta=\pi$. Or un thyristor ne peut conduire un courant n\'egatif : il se bloque exactement au passage par z\'ero du courant, soit $\theta = \pi$ (blocage naturel, comme pour une diode).
  \item Sur une p\'eriode $2\pi$, $u_S(\theta) = V_m\sin\theta$ pour $\theta\in[\alpha\,;\,\pi]$ et $u_S(\theta)=0$ ailleurs :
  \[
  \langle u_S \rangle = \dfrac{1}{2\pi}\int_{\alpha}^{\pi} V_m\sin\theta\,d\theta
  = \dfrac{V_m}{2\pi}\Big[-\cos\theta\Big]_{\alpha}^{\pi}
  = \dfrac{V_m}{2\pi}\big(-\cos\pi + \cos\alpha\big)
  = \dfrac{V_m}{2\pi}(1+\cos\alpha)
  \]
  \item Pour $\alpha = \SI{60}{\degree}$ : $\langle u_S\rangle = \dfrac{325}{2\pi}(1+\cos 60^\circ) = 51{,}7\times(1+0{,}5) \approx \SI{77,6}{\volt}$.

  Pour $\alpha = 0$ (\'equivalent \`a une diode non command\'ee) : $\langle u_S\rangle = \dfrac{325}{2\pi}(1+1) = \dfrac{325}{\pi} \approx \SI{103,5}{\volt}$, qui est bien la valeur moyenne classique d'un redressement demi-onde \`a diode. Retarder l'amor\c cage ($\alpha$ croissant) r\'eduit donc la dur\'ee de conduction et, par cons\'equent, la valeur moyenne d\'elivr\'ee \`a la charge~: l'angle d'amor\c cage $\alpha$ est le param\`etre de r\'eglage continu de $\langle u_S\rangle$ entre $V_m/\pi$ (pour $\alpha=0$) et $0$ (pour $\alpha=\pi$).
\end{enumerate}

\textbf{Conclusion.} Le montage P1 command\'e r\`egle donc la tension moyenne de charge de \SI{103,5}{\volt} \`a $0$ sans aucun organe dissipatif, mais au prix d'une ondulation tr\`es forte (conduction sur moins d'une demi-p\'eriode) : ce crit\`ere de qualit\'e de l'onde interdit son emploi d\`es qu'une tension continue liss\'ee est exig\'ee, et impose alors un montage PD2 ou PD3 command\'e associ\'e \`a un filtrage.
}
\end{tcolorbox}

\CAPETENDFILE
